\section{Quantum divisors and point descendants}\label{sec:setup}

We set up a finite collection of commuting matrices and a sequence of
vectors on which they act. Their grading records anticanonical degrees;
the vectors will connect their joint eigenvalues to rational curves.
Throughout the proof, $X$ is a smooth projective complex Fano manifold,
$n=\dim X>0$, $r=\rho_X$, and $\iota=\iota_X\ge2$.
The case $\iota=1$ will require no argument beyond $n>0$.

Choose an integral basis $D_1,\ldots,D_r$ of the N\'eron--Severi group
modulo torsion. A numerical curve degree is represented by
\[
 \beta=(D_1\cdot\beta,\ldots,D_r\cdot\beta)\in\mathbb Z^r.
\]
Here and below we use the same symbol for a numerical degree and its
coordinate vector. Write
\[
 -K_X\equiv\sum_{j=1}^r c_jD_j,
 \qquad d_\beta=\sum_{j=1}^r c_j\beta_j.
\]
A positive degree means the degree of a nonconstant genus-zero stable
map. Every such degree has $d_\beta\ge\iota$. There are only finitely
many positive numerical degrees with bounded $d_\beta$: for a
sufficiently large integer $a$, both $a(-K_X)+D_j$ and $a(-K_X)-D_j$
are ample, so $|\beta_j|\le a d_\beta$ on effective degrees.
When several curve classes have the same numerical degree, we sum their
invariants. Proper finite-type spaces of stable maps of bounded
projective degree make these sums finite.

We use ordinary genus-zero Gromov--Witten invariants, defined by the
virtual fundamental classes of stable-map stacks. An insertion
$\tau_\ell(a)$ means $\psi^\ell\ev^*a$ at the indicated marking:
$\ev$ evaluates the stable map there, and $\psi$ is the first Chern
class of its cotangent line.
The divisor, string and topological recursion equations, together with
associativity, are the standard identities we use
\cite{BehrendVirtual,KontsevichManinQuantum,KontsevichManinDescendants,PandharipandeDescendants}.
The discussion below specifies all unstable conventions needed here.
Set
\[
 H=\bigoplus_{k=0}^n H_k,
 \qquad H_k=H^{2k}(X,\mathbb C),
 \qquad (a,b)=\int_Xa\cup b,
\]
and put $H_k=0$ outside this range. The pairing on $H$ is nondegenerate.
We work only with even cohomology; invariants involving a single odd
intermediate class and otherwise even insertions vanish by parity.

For a positive numerical degree $\beta$, define $S_\beta\in\operatorname{End}(H)$ by
\[
 (S_\beta a,b)=\langle a,b\rangle_\beta.
\]
The virtual dimension of the two-pointed space is $n+d_\beta-1$.
Consequently
\begin{equation}\label{eq:grading}
 S_\beta(H_k)\subset H_{k+1-d_\beta}.
\end{equation}
In particular $S_\beta=0$ for $d_\beta>n+1$. For nonzero formal variables
$q_1,\ldots,q_r$, write $q^\beta=\prod_jq_j^{\beta_j}$ and define
\begin{equation}\label{eq:quantum-operators}
 A_j(q)=D_j\cup(-)+\sum_{\beta>0}q^\beta\beta_jS_\beta
       =\sum_\gamma q^\gamma A_{j,\gamma}.
\end{equation}
Thus $A_{j,0}=D_j\cup(-)$, $A_{j,\beta}=\beta_jS_\beta$ for positive
degrees, and all other coefficients vanish. The grading bound and
finiteness of bounded degrees make these Laurent polynomial matrices.
The divisor equation identifies them with small quantum multiplication
by $D_j$, so associativity implies that they commute. If $R_t$ acts
by $t^k$ on $H_k$, then
\begin{equation}\label{eq:homogeneity}
 A_j(t^{c_1}q_1,\ldots,t^{c_r}q_r)
       =tR_t^{-1}A_j(q)R_t.
\end{equation}
Indeed the classical term raises the grading by one, whereas the
$\beta$ term changes it by $1-d_\beta$.

\subsection{The descendant recurrence}

Two-point descendant generating series also occur in Givental's
quantum differential framework~\cite[\S6]{Givental1996}.  The
recurrence below is derived directly from the virtual genus-zero
identities cited above, so it does not require a convexity assumption
on $X$.

Let $\pt\in H_n$ be the class of a point, normalized by
$\int_X\pt=1$. Define $v_0=\pt$ and, for positive degrees, define
$v_\beta\in H$ by
\[
 (v_\beta,b)=\sum_{\ell\ge0}
             \langle\tau_\ell(\pt),b\rangle_\beta.
\]
The dimension constraint makes each sum finite. Set all other
$v_\beta$ equal to zero. In particular there is no degree-zero
unstable two-point invariant in this definition.

\begin{lemma}\label{lem:recurrence}
For every integer degree vector $\beta$ with $d_\beta>0$ and every
$j\in\{1,\ldots,r\}$,
\begin{equation}\label{eq:recurrence}
 \beta_jv_\beta=\sum_\gamma A_{j,\gamma}v_{\beta-\gamma}.
\end{equation}
For a positive degree $\beta$,
\begin{equation}\label{eq:string-pairing}
 (v_\beta,1)=\langle\tau_{d_\beta-2}(\pt)\rangle_\beta.
\end{equation}
\end{lemma}

\begin{proof}
Pair the first identity with an arbitrary even class $b$. Inserting
$D_j$ into $\langle\tau_\ell(\pt),b\rangle_\beta$ multiplies the invariant
by $\beta_j$ \cite[Lemma~1.4, equation~(12)]{KontsevichManinDescendants}: the descendant correction contains
$D_j\cup\pt=0$. For $\ell>0$, genus-zero topological recursion at the
point marking \cite[equation~(4a)]{KontsevichManinDescendants} separates a two-point factor containing
$\tau_{\ell-1}(\pt)$ from a three-point factor containing $D_j$ and $b$.
The point side has positive degree, since a two-point degree-zero
component is unstable. The other side can have degree zero, in which
case it contributes the classical cup operator. Summing over $\ell>0$
therefore gives the terms with positive-degree $v_{\beta-\gamma}$.
The $\ell=0$ primary invariant gives the remaining term
$A_{j,\beta}v_0$. Splitting via the Poincar\'e pairing gives exactly
\eqref{eq:recurrence}. If a degree is not represented, the same splitting
identity holds with zero coefficients.

For the second identity, dimension leaves only $\ell=d_\beta-1$ in the
pairing with $1$. The string equation \cite[\S1.2, p.~311]{PandharipandeDescendants} then gives
\[
 \langle\tau_{d_\beta-1}(\pt),1\rangle_\beta
 =\langle\tau_{d_\beta-2}(\pt)\rangle_\beta.
\]
The exponent is nonnegative because $d_\beta\ge\iota\ge2$.
\end{proof}

The next section gives a positive lower bound for this pairing whenever
$\beta$ has a free rational representative. It is a statement about a
pushforward of the virtual evaluation fiber; it will not require the
entire virtual class to be effective.
